\section{Discussion}
\label{sec:discussion}

The results attribute the estimator's behavior to its design choices with a
regularity worth stating plainly. \emph{Weighting is the dominant choice
wherever the design is informative}: on the \scf{}, whose list sample
oversamples wealthy households, removing the weighted bootstrap costs a
factor of six on net-worth marginal fit and a factor of 2.3 on the
population-view energy distance --- larger than the gap between the candidate
and any competing method. On the \cps{} components task, whose derived person
weights carry little outcome information, weighting is approximately free.
An estimator cannot know in advance which regime it is in, which is the
argument for weighting by construction rather than by option. \emph{Gating
carries the zero-inflation surface} (0.6 versus 4.5 percentage points of
dividend zero-share error against the same forest ungated), and
\emph{chaining} contributes a consistent, smaller improvement on dependent
targets (17 percent of net-worth $W_1$). None of the three hurts where it
does not help.

The evaluation findings are as load-bearing as the method findings. Marginal
metrics saturate: with weak predictors, an unconditional weighted draw from
the donor marginal ties or beats every conditional method on pinball loss and
Wasserstein distance --- on the OpenML suite it attains the best mean rank
--- while carrying no conditional structure at all. The population-view
harness is what separates it: the classifier two-sample test detects its
broken demographics--wealth coupling (AUC 0.903 against the candidate's
0.843), and the coverage/energy decomposition localizes \emph{which}
ingredient of a candidate fails --- support or measure --- with direct
consequences for whether downstream calibration can repair it. The fragility
diagnostic adds the complementary caution: the goal is the truth's own tail
exposure, not minimal exposure, and the one method that minimizes it (OLS)
does so by failing to generate realistic tails at all.

Several limitations bound the claims. First, the population-view experiment
runs one view pair: the common \cps{}--\scf{} offset (differing income and
demographic measurement, the dropped education predictor) is shared by all
methods but bounds how close any candidate can come to the floor; a
multi-view instantiation (adding \sipp{} and \psid{}) would tighten the
identification and is the production release-gate use of the harness.
Second, the pairwise joint metrics run on weight-proportional subsamples
(2,048 points per side), which bounds their resolution; the C2ST, which uses
all rows, is the axis least affected. Third, baseline quantile models are
converted to samplers through a 99-level quantile grid, which truncates
draws beyond the 1st and 99th conditional percentiles; populace-fit samples
continuous quantiles, so extreme-tail comparisons slightly favor it by
construction, though the affected mass is under 2 percent. Fourth, the
\microimpute{} forest's internal randomness is not seedable through its
public constructor, so its paired-seed dispersion includes forest noise the
candidate's does not. Fifth, comparisons hold shared forest backends and
default hyperparameters fixed rather than tuning each method per task; the
ablations are exact controls, the cross-family comparisons are
defaults-versus-defaults. Finally, we evaluate single imputations
distributionally and do not propagate imputation uncertainty in the
multiple-imputation sense \citep{rubin1987multiple}; for the population
statistics microsimulation consumes, the distributional view is the relevant
one, but inference from imputed files would require the fuller treatment.